% Pré-projeto HIJACK — arquivo LaTeX editável
\documentclass[12pt,a4paper]{article}

% ---------- Idioma e codificação ----------
\usepackage[utf8]{inputenc}
\usepackage[T1]{fontenc}
\usepackage[brazil]{babel}
\usepackage{lmodern}

% ---------- Layout (próximo da ABNT) ----------
\usepackage[top=3cm,left=3cm,bottom=2cm,right=2cm]{geometry}
\usepackage{setspace}
\onehalfspacing
\setlength{\parindent}{1.25cm}
\setlength{\parskip}{0.2cm}

% ---------- Pacotes gerais ----------
\usepackage{graphicx}
\usepackage{tikz}
\usetikzlibrary{calc}
\usepackage{float}
\usepackage{booktabs}
\usepackage{tabularx}
\usepackage{array}
\usepackage{enumitem}
\usepackage[table]{xcolor}
\usepackage{titlesec}
\usepackage[hidelinks]{hyperref}
\usepackage{xurl}

% ---------- Estilo dos títulos ----------
\titleformat{\section}{\normalfont\large\bfseries\MakeUppercase}{\thesection}{1em}{}
\titleformat{\subsection}{\normalfont\normalsize\bfseries}{\thesubsection}{1em}{}
\titleformat{\subsubsection}{\normalfont\normalsize\itshape}{\thesubsubsection}{1em}{}

% ---------- Metadados do documento (EDITAR) ----------
\newcommand{\instituicao}{Universidade Federal do Rio Grande do Norte (UFRN)}
\newcommand{\disciplina}{Fundamentos de Circuitos e Sistemas Controlados}
\newcommand{\titulo}{HIJACK: demonstração de segurança em carrinhos controlados remotamente}
\newcommand{\orientador}{Alexandre Magnus Fernandes Guimarães}
\newcommand{\cidade}{Natal -- RN}
\newcommand{\ano}{2026}


\begin{document}
\begin{titlepage}
\centering
{\large\bfseries \instituicao}\\[0.3cm]
{\large \disciplina}
\vfill
{\normalsize
Matheus Felipe Duarte Dionísio\\
João Gonçalves de Souza Neto\\
Élisson Rian Félix Simas\\
Carlos Roberto de Miranda Gomes Neto\\
Gustavo Henrique Assunção Paiva\\
Pedro Nascimento de Paiva Fernandes Neto\\
Marcelo Jorge Oliveira de Albuquerque Melo}
\vfill
{\Large\bfseries \MakeUppercase{\titulo}}\\[0.6cm]
{\large Pré-projeto}
\vfill
\begin{flushright}
\begin{minipage}{0.55\textwidth}
\singlespacing
Pré-projeto apresentado à disciplina de \disciplina{},
na \instituicao.\\[0.3cm]
Professor: \orientador
\end{minipage}
\end{flushright}
\vfill
{\large \cidade\\ \ano}
\end{titlepage}

\section*{Resumo}
\addcontentsline{toc}{section}{Resumo}
Este pré-projeto apresenta o desenvolvimento do HIJACK, uma atividade prática de
cibersegurança que utiliza carrinhos controlados por uma rede
local. A proposta é dividir os participantes em duas equipes, cada uma com um
piloto e um responsável pelos ataques e pelas defesas. Enquanto o piloto
controla o carrinho pelo celular, o outro integrante utiliza um painel no
notebook para tentar interferir nos comandos do adversário e proteger os da
própria equipe. A demonstração começa com um sistema intencionalmente
vulnerável: uma alteração na interpretação dos comandos pode fazer o carrinho
andar para trás quando o piloto tenta avançar, por exemplo. Depois, a equipe demonstrará a troca da senha padrão por uma senha mais forte e o bloqueio simulado das alterações de comando.
A intenção é mostrar, por meio do comportamento físico dos carrinhos, por que
não basta que um dispositivo conectado receba comandos: ele também precisa
verificar sua origem e sua validade. O trabalho relaciona essa experiência a
incidentes de segurança em sistemas conectados e ao uso recente de inteligência
artificial na identificação e exploração de vulnerabilidades.

\noindent\textbf{Palavras-chave:} cibersegurança; sistemas ciberfísicos;
controle remoto; integridade de comandos; ensino prático.
\newpage
\tableofcontents
\newpage

\section{Introdução}
É comum associar ataques cibernéticos ao roubo de informações ou à interrupção
de serviços digitais. Quando um programa controla um equipamento físico, porém,
os efeitos podem aparecer fora da tela. Um comando recebido ou interpretado de
forma incorreta pode mudar o comportamento de um robô, de um veículo ou de
outro dispositivo conectado.

Em 2015, pesquisadores demonstraram que era possível interferir remotamente
em funções de um Jeep Cherokee explorando falhas no sistema conectado do
veículo. O episódio levou a um recall de aproximadamente 1,4 milhão de
unidades \cite{wired2015jeep}. Esse caso ajuda a entender por que autenticação,
controle de acesso e integridade de mensagens também são relevantes em
sistemas que atuam no mundo físico.

A ideia do projeto surgiu da possibilidade de transformar essa discussão em
uma demonstração observável. Em vez de apenas apresentar um exemplo de
comando adulterado, queremos que o público veja o efeito no carrinho. Se o
piloto mandar seguir em frente e o veículo recuar, fica mais fácil discutir
onde a comunicação falhou e o que poderia ter evitado o problema.

O HIJACK será desenvolvido a partir de um protótipo simples, com motores DC,
microcontrolador e interface web. A proposta inclui uma etapa vulnerável e
outra com medidas de proteção, para permitir a comparação entre os dois
comportamentos.

A montagem também permite relacionar os conceitos de circuitos e sistemas
controlados com situações práticas. O acionamento dos motores pela ponte H, o posicionamento da direção por
servo motor, a alimentação do ESP32 e o envio de sinais mostram como decisões
de software interferem diretamente no funcionamento de um circuito elétrico.

\section{Fundamentação e Justificativa}
\subsection{O tempo disponível para corrigir falhas}
A divulgação de uma vulnerabilidade nem sempre é seguida imediatamente pela
correção dos sistemas afetados. Entre a descoberta da falha, a publicação de
uma atualização e sua instalação, pode existir um período em que os usuários
continuam expostos. Relatórios publicados em 2026 destacam a rapidez com que
algumas falhas passaram a ser exploradas \cite{synack2026,thn2026bas}.

Uma análise divulgada pela CrowdStrike também aponta exploração de falhas
pouco tempo depois de surgirem provas de conceito públicas
\cite{crowdstrike2026}. Em outro exemplo, a Cloud Security Alliance descreveu
um caso envolvendo o framework PraisonAI no qual uma exploração funcional foi
produzida em poucas horas \cite{csa2026praison}. Os números variam conforme o
levantamento e o tipo de vulnerabilidade, mas apontam para a mesma preocupação:
não é prudente depender apenas da aplicação posterior de correções.

No carrinho, esse problema será representado por um sistema que inicialmente
aceita comandos sem checar adequadamente quem os enviou. A atividade permitirá
discutir o impacto dessa decisão e avaliar como o comportamento muda depois da
inclusão de verificações.

\subsection{Inteligência artificial em ataques cibernéticos}
O uso de modelos de IA em atividades ofensivas tornou a discussão sobre
segurança ainda mais atual. Em novembro de 2025, a Anthropic relatou uma
campanha identificada como GTG-1002, na qual agentes de IA foram empregados
para executar uma parcela significativa das etapas operacionais de uma
intrusão, enquanto humanos continuavam responsáveis por decisões importantes
\cite{thn2025gtg1002}.

Segundo os relatos publicados sobre o caso, a automação abrangeu tarefas como
reconhecimento de alvos, investigação de falhas e coleta de informações
\cite{bdtech2025,aimag2025}. Em 2026, a Check Point também discutiu o uso de IA
na execução de etapas de ataques \cite{checkpoint2026}. Isso não significa que
todo ataque seja autônomo, mas mostra como ferramentas de automação podem
reduzir o trabalho necessário para certas atividades.

A dinâmica proposta não depende de utilizar IA para invadir os carrinhos.
A relação com esses casos está na sequência das ações: identificar o alvo,
localizar uma fraqueza, tentar explorá-la e observar se a proteção existente
é suficiente.

\subsection{O comprometimento do LiteLLM}
Outro aspecto importante é a confiança em componentes de software. Em março
de 2026, versões distribuídas do pacote LiteLLM no PyPI foram comprometidas,
permitindo a execução de código malicioso e a tentativa de obtenção de
credenciais \cite{arthur2026litellm,infoq2026litellm}.

O caso foi associado ao comprometimento anterior de uma ferramenta utilizada
em processos de desenvolvimento, o Trivy. A partir desse acesso, os atacantes
conseguiram interferir na publicação do pacote \cite{netenrich2026litellm,
hackernoon2026litellm}. Algumas análises destacaram ainda o uso de um arquivo
\texttt{.pth}, capaz de executar instruções durante a inicialização do Python
\cite{bastion2026litellm}.

Esse incidente interessa ao projeto porque evidencia um problema de confiança:
mesmo quando o usuário acredita estar utilizando um componente legítimo, algo
no caminho pode ter sido alterado. No HIJACK, essa ideia será apresentada de
forma simplificada pela adulteração da relação entre um comando recebido e
a ação executada pelo motor. Não se trata de reproduzir o ataque ao LiteLLM,
mas de usar o caso como referência para explicar a importância da integridade.

\subsection{Inteligência artificial como ferramenta de defesa}
A IA também tem sido utilizada na análise de código e na descoberta de
vulnerabilidades. Entre os exemplos recentes estão as falhas apelidadas de
PHP2Shell, divulgadas em 2026, identificadas com auxílio de ferramentas de IA
\cite{wiz2026php2shell,fortune2026php}. Outro caso é o Node2Shell, uma
vulnerabilidade relacionada ao processamento Brotli no Node.js, para a qual
foram publicados alertas e correções de segurança
\cite{socprime2026node,akamai2026node2shell}.

Há ainda experiências mais amplas, como o AI Cyber Challenge, voltado ao uso
de sistemas automatizados para encontrar e corrigir problemas de segurança
em software \cite{drishti2026}. Esses exemplos reforçam que ferramentas de IA
podem ser úteis nos dois lados do problema: tanto na identificação de pontos
fracos quanto no trabalho de corrigi-los.

No nosso projeto, o foco principal não será construir um agente de IA. O
objetivo é demonstrar um princípio de segurança que continua válido
independentemente de quem realiza o ataque: um dispositivo não deve confiar
em comandos apenas porque eles chegaram ao seu canal de comunicação.

\subsection{Relação dos casos com a proposta}
\begin{table}[H]
\centering\small\renewcommand{\arraystretch}{1.3}
\begin{tabularx}{\textwidth}{>{\raggedright\arraybackslash}X >{\raggedright\arraybackslash}X}
\toprule
\textbf{Situação discutida} & \textbf{Exemplo na demonstração} \\
\midrule
Controle remoto de equipamentos & O movimento do carrinho muda após um comando indevido \\
Confiança em componentes intermediários & A interpretação dos comandos pode ser modificada \\
Reconhecimento e exploração de falhas & O painel do hacker organiza as tentativas de ataque por etapas \\
Verificação de integridade e origem & A plataforma simula o bloqueio de adulterações e autentica o acesso por senha \\
\bottomrule
\end{tabularx}
\caption{Relação entre os conceitos estudados e a demonstração do HIJACK.}
\label{tab:relacao}
\end{table}

\section{Objetivos}
\subsection{Objetivo geral}
Construir e apresentar um sistema de carrinhos controlados remotamente que
permita demonstrar, em um ambiente de testes, os efeitos de comandos
adulterados e a importância de mecanismos de proteção na comunicação com
dispositivos físicos.

\subsection{Objetivos específicos}
\begin{itemize}[leftmargin=*]
\item Adaptar a estrutura do carrinho com motores DC para tração e um servo motor para esterçamento das rodas;
\item Programar o microcontrolador para receber e executar os comandos enviados pela interface de controle;
\item Desenvolver uma aplicação web com telas específicas para o piloto e para o responsável pelo ataque e defesa;
\item Criar uma versão propositalmente vulnerável, restrita ao ambiente de testes do grupo;
\item Demonstrar a alteração de comandos e observar seus efeitos no movimento do carrinho;
\item Implementar autenticação por senha e demonstrar, de forma simulada, as proteções de integridade e contra repetição;
\item Realizar uma demonstração interativa e livre, com duas equipes atuando simultaneamente, e avaliar o comportamento do sistema.
\end{itemize}

\section{Descrição do Projeto}
\subsection{Funcionamento geral}
O projeto recebeu o nome \textbf{HIJACK}. Ele será composto por dois
carrinhos controlados por placas ESP32 e por uma aplicação web executada
localmente em um notebook. Os celulares dos pilotos, os notebooks dos
participantes e os microcontroladores serão conectados à mesma rede Wi-Fi.
A troca de comandos entre o servidor e os ESP32 ocorrerá por WebSocket.

Cada equipe contará com um piloto, que dirige seu carrinho pelo celular,
e um hacker, que usa o painel web para interferir no carrinho adversário e
proteger o próprio equipamento. As duas equipes atuarão simultaneamente.
Não haverá pontuação, classificação ou vencedores; a atividade ficará
disponível enquanto houver participantes interessados.

A Figura~\ref{fig:arquitetura} apresenta o fluxo previsto de comunicação.
\begin{figure}[H]
\centering
\begin{tikzpicture}[font=\small,>=stealth]
\node[draw,rounded corners,align=center,minimum width=3.7cm,minimum height=1cm] (piloto) at (-3,2) {Celular do piloto\\Interface de controle};
\node[draw,rounded corners,align=center,minimum width=3.7cm,minimum height=1cm] (hacker) at (3,2) {Notebook do hacker\\Painel de ataque e defesa};
\node[draw,rounded corners,align=center,minimum width=4.6cm,minimum height=1cm] (servidor) at (0,0) {Servidor web HIJACK\\Notebook na rede local};
\node[draw,rounded corners,align=center,minimum width=3.8cm,minimum height=1cm] (a) at (-3,-2) {ESP32 + L298N\\Carrinho A};
\node[draw,rounded corners,align=center,minimum width=3.8cm,minimum height=1cm] (b) at (3,-2) {ESP32 + L298N\\Carrinho B};
\draw[<->] (piloto.south) -- (servidor.north west);
\draw[<->] (hacker.south) -- (servidor.north east);
\draw[<->] (servidor.south west) -- node[left,font=\scriptsize] {WebSocket} (a.north);
\draw[<->] (servidor.south east) -- node[right,font=\scriptsize] {WebSocket} (b.north);
\end{tikzpicture}
\caption{Arquitetura de comunicação prevista para a plataforma HIJACK.}
\label{fig:arquitetura}
\small Fonte: elaborado pelos autores (2026).
\end{figure}

\subsection{Montagem do protótipo}
A primeira montagem foi feita com uma base de papelão, utilizada para prender
os motores, as rodas e os demais componentes. Essa estrutura é provisória e será substituída por um chassi fabricado
por impressão 3D, para facilitar a fixação e a organização dos componentes. A Figura~\ref{fig:carrinho} apresenta o estado atual do protótipo.

\begin{figure}[H]
\centering
\includegraphics[height=0.53\textheight]{figuras/prototipo_carrinho.jpg}
\caption{Montagem inicial do carrinho sobre uma base provisória de papelão.}
\label{fig:carrinho}
\small Fonte: acervo do grupo (2026).
\end{figure}

A montagem inicial utiliza dois motores DC com redução, uma ponte H para controlar o
sentido de rotação, fios de ligação e um suporte para baterias. Na versão
planejada, será acrescentado um servo motor com ligação mecânica ao conjunto
de direção, conforme o esquema apresentado pelo grupo. O grupo informou que as ligações dos GPIOs já foram conferidas. O
esquema elétrico é apresentado como preliminar porque a estabilidade da
alimentação ainda precisa ser testada sob carga.

\begin{table}[H]
\centering\small\renewcommand{\arraystretch}{1.3}
\begin{tabularx}{\textwidth}{l >{\raggedright\arraybackslash}X}
\toprule
\textbf{Componente} & \textbf{Descrição} \\
\midrule
Microcontrolador & ESP32, responsável pela comunicação via Wi-Fi, tração e direção \\
Controle dos motores & Ponte H L298N \\
Tração & Dois motores DC com caixa de redução e rodas \\ 
Direção & Servo motor e mecanismo de ligação às rodas direcionais (previstos) \\
Alimentação & Duas células 18650 de 3,7\,V, ligadas em série (7,4\,V nominais) \\
Estrutura & Base provisória de papelão; chassi definitivo impresso em 3D \\
Comunicação & Wi-Fi em rede local, com troca de mensagens via WebSocket \\
\bottomrule
\end{tabularx}
\caption{Componentes previstos e utilizados na montagem inicial.}
\label{tab:componentes}
\end{table}

\subsection{Circuito, tração e direção}
O ESP32 enviará sinais digitais à ponte H L298N, que controlará o sentido
de rotação dos dois motores DC responsáveis pela tração. A velocidade
será fixa. O movimento para frente e para trás dependerá do sentido
selecionado na ponte H.

Para realizar curvas, o projeto prevê uma direção por servo motor.
O servo será acoplado mecanicamente ao conjunto de rodas direcionais,
alterando seu ângulo em relação ao chassi. Assim, ao solicitar uma curva,
o ESP32 comandará a posição do servo, enquanto os motores DC continuarão
responsáveis pelo deslocamento. Ao contrário da proposta inicial de
tração diferencial, as curvas não serão obtidas fazendo os motores
girarem em sentidos opostos. O desenho mecânico ainda será detalhado
na montagem do chassi impresso em 3D.

O conjunto será alimentado por duas baterias 18650 de 3,7\,V em série,
com tensão nominal de 7,4\,V e tensão máxima de aproximadamente 8,4\,V
quando completamente carregadas. Na montagem informada, o ESP32 recebe
alimentação da saída de 5\,V do módulo L298N. Essa ligação exige validação:
a disponibilidade de 5\,V regulados pode variar conforme o módulo e a
tensão de entrada, e picos de corrente dos motores podem causar quedas
de tensão. Se necessário, será utilizado um regulador separado para
alimentar o microcontrolador com segurança. A alimentação do servo motor
também deverá ser dimensionada e verificada, considerando sua corrente
de partida e de travamento. O servo não deve ser alimentado diretamente
por um GPIO do ESP32. Todas as partes devem compartilhar uma referência
de terra comum.

\begin{figure}[H]
\centering
\begin{tikzpicture}[font=\small,>=stealth,
  every node/.style={align=center}]
\node[draw,rounded corners,minimum width=3cm,minimum height=0.9cm] (bat) at (0,2.7) {2 baterias 18650\\7,4\,V nominais};
\node[draw,rounded corners,minimum width=2.7cm,minimum height=0.9cm] (ponte) at (0,0.7) {Ponte H L298N};
\node[draw,rounded corners,minimum width=2.5cm,minimum height=0.9cm] (esp) at (-4.1,0.7) {ESP32\\Wi-Fi / WebSocket};
\node[draw,rounded corners,minimum width=1.7cm,minimum height=0.75cm] (ma) at (3.8,1.5) {Motor A};
\node[draw,rounded corners,minimum width=1.7cm,minimum height=0.75cm] (mb) at (3.8,-0.1) {Motor B};
\node[draw,rounded corners,minimum width=2.2cm,minimum height=0.75cm] (servo) at (-4.1,-1.7) {Servo motor\\(direção)};
\node[draw,rounded corners,minimum width=3cm,minimum height=0.75cm] (rodas) at (0,-2.1) {Mecanismo de direção\\e rodas direcionais};
\draw[->] (bat) -- node[right,font=\scriptsize] {via chave} (ponte);
\draw[->] (esp) -- node[above,font=\scriptsize] {sinais de controle} (ponte);
\draw[->] (ponte.east) -- ++(.55,0) |- (ma.west);
\draw[->] (ponte.east) -- ++(.55,0) |- (mb.west);
\draw[->] (esp.south) -- node[left,font=\scriptsize] {sinal PWM de posição} (servo.north);
\draw[->] (servo.east) -- (rodas.west);
\node[font=\scriptsize] at (-1.95,-2.78) {ligação mecânica};
\draw[->,dashed] (bat.west) -| node[pos=.85,left,font=\scriptsize] {alimentação\\a validar} (esp.north);
\node[font=\scriptsize] at (0,-3.55) {GND comum; alimentação do servo e do ESP32 a verificar};
\end{tikzpicture}
\caption{Diagrama funcional preliminar da tração e da direção por servo motor. Não representa a pinagem definitiva nem a disposição mecânica das rodas.}
\label{fig:circuito}
\small Fonte: elaborado pelos autores (2026).
\end{figure}

O grupo informou que os GPIOs da montagem de tração já foram conferidos.
A tabela registra essas ligações; o pino de sinal do novo servo motor
ainda será definido e testado, assim como sua alimentação. O diagrama
funcional não substitui a conferência física da montagem.

\begin{table}[H]
\centering\small\renewcommand{\arraystretch}{1.25}
\begin{tabular}{l l l}
\toprule
\textbf{Sinal do ESP32} & \textbf{Ligação na L298N} & \textbf{Finalidade} \\
\midrule
GPIO14 & ENA & Habilitação do motor A \\
GPIO27 / GPIO26 & IN1 / IN2 & Sentido de rotação do motor A \\
GPIO25 / GPIO33 & IN3 / IN4 & Sentido de rotação do motor B \\
GPIO32 & ENB & Habilitação do motor B \\
GND & GND & Referência comum \\
\bottomrule
\end{tabular}
\caption{Ligações de sinais informadas pelo grupo.}
\label{tab:ligacoes}
\end{table}

\subsubsection{Comandos de movimento}
A lógica de controle será dividida entre tração e direção. Os comandos
\texttt{FRENTE} e \texttt{RE} definirão o sentido de rotação dos motores DC
pela ponte H. Já os comandos \texttt{ESQUERDA} e \texttt{DIREITA} ajustarão
a posição do servo motor. Na posição central, as rodas permanecerão
alinhadas para seguir em linha reta. Os ângulos efetivos do servo serão
ajustados durante a montagem, conforme os limites mecânicos do conjunto.

\begin{table}[H]
\centering\small\renewcommand{\arraystretch}{1.3}
\begin{tabularx}{\textwidth}{l >{\raggedright\arraybackslash}X}
\toprule
\textbf{Comando} & \textbf{Resposta prevista} \\
\midrule
Frente & Acionar os motores DC para avançar \\
Ré & Acionar os motores DC para recuar \\
Esquerda & Posicionar o servo para esterçar as rodas à esquerda \\
Direita & Posicionar o servo para esterçar as rodas à direita \\
Centro & Retornar o servo à posição central \\
Parar & Desligar a tração; centralização da direção conforme a lógica do controle \\
\bottomrule
\end{tabularx}
\caption{Funções previstas para tração e direção, sujeitas aos testes do protótipo.}
\label{tab:movimentos}
\end{table}

Os comandos de deslocamento e direção poderão ser combinados durante o uso
do analógico virtual. A adulteração simulada na plataforma poderá trocar
frente por ré ou esquerda por direita. No segundo caso, o efeito observado
será o esterçamento para o lado oposto ao solicitado, e não a inversão do
sentido de rotação dos motores de tração.

\subsection{Interfaces web}
A plataforma terá uma tela de entrada e dois painéis separados. Essa divisão
foi escolhida porque o piloto precisa de controles simples e rápidos,
enquanto o hacker precisa acompanhar as condições da partida e as opções de
ataque e defesa.

\subsubsection{Entrada na partida}
Na tela inicial, o participante informa o nome, escolhe a função e seleciona
o carrinho da equipe. Também há um campo de senha, que poderá fazer parte da
demonstração sobre autenticação. No modo simulado, a aplicação utiliza a senha
\texttt{1234}; isso não representa uma configuração segura para uso real.

\begin{figure}[H]
\centering
\includegraphics[width=0.60\textwidth]{figuras/tela_login.jpg}
\caption{Tela de entrada da aplicação HIJACK.}
\label{fig:login}
\small Fonte: acervo do grupo (2026).
\end{figure}

\subsubsection{Painel do piloto}
O painel do piloto apresenta um controle analógico virtual e dois botões de
ação, identificados como A e B. A interface também exibe um estado de espera antes do início da demonstração. Os botões A e B aparecem como elementos visuais da interface,
sem função de controle definida nesta etapa. A velocidade dos motores
será fixa e o carrinho parará quando o piloto soltar o controle.

\begin{figure}[H]
\centering
\includegraphics[width=0.75\textwidth]{figuras/controle_piloto.jpg}
\caption{Interface de controle utilizada pelo piloto.}
\label{fig:piloto}
\small Fonte: acervo do grupo (2026).
\end{figure}

\subsubsection{Painel do hacker}
O painel do hacker possui abas de Ataque e Defesa, informações de conexão
e um indicador visual dos mecanismos de proteção. O ataque será simulado
pelo servidor, com etapas guiadas: escanear a rede, identificar o carrinho
adversário, explorar uma falha simulada e modificar os comandos.

Depois de concluir essas etapas, o participante poderá escolher quais
movimentos deseja inverter: frente e ré, esquerda e direita. A inversão
será aplicada pelo próprio servidor antes de encaminhar os comandos ao
ESP32. Não haverá interceptação real de tráfego nem exploração de redes
ou equipamentos externos.

\begin{figure}[H]
\centering
\includegraphics[width=\textwidth]{figuras/painel_hacker.jpg}
\caption{Painel do hacker antes do início da demonstração.}
\label{fig:hacker}
\small Fonte: acervo do grupo (2026).
\end{figure}

\subsection{Dinâmica da apresentação}
A atividade será realizada simultaneamente por duas equipes. Cada piloto
controlará seu carrinho pelo celular enquanto seu colega utilizará o painel
para atacar o equipamento adversário e defender o próprio. A interação
ocorrerá livremente, sem duração fixa, pontuação, corrida ou vencedores.

O público poderá observar a resposta do carrinho às mudanças simuladas
de comando e, em seguida, o bloqueio dessas mudanças após a ativação da
defesa correspondente. A apresentação se concentrará na diferença entre
o comando solicitado e o executado pelos motores.

\subsection{Vulnerabilidade planejada}
O sistema de testes começará com a senha padrão \texttt{1234} para ambos
os carrinhos. Essa credencial fraca será utilizada propositalmente para
discutir os riscos de senhas previsíveis. O servidor também oferecerá
uma simulação de adulteração dos comandos: quando a opção estiver ativa,
poderá trocar individualmente frente por ré ou esquerda por direita.
A simulação não pressupõe que o participante tenha comprometido a rede.

\subsection{Medidas de proteção}
O mecanismo real de autenticação será demonstrado por meio da substituição
da senha padrão por uma senha mais forte. O servidor passará a exigir a
nova credencial para permitir o acesso ao controle. Isso ilustra a
importância de impedir o uso de senhas padrão, mas não impede por si só
alterações nos comandos feitas por um componente autorizado.

As camadas de integridade e proteção contra repetição serão apresentadas
como \textbf{simulações didáticas} na interface. Ao ativar a defesa de
integridade, o servidor deixará de aplicar as inversões solicitadas pelo
hacker. A proteção contra repetição exibirá o conceito de recusar comandos
antigos, sem implementar ainda validação criptográfica ou contadores no
firmware. A adoção de mecanismos como HMAC e números de sequência poderá
ser considerada em versões futuras.

Como proteção física, o ESP32 deverá interromper os motores quando o
piloto soltar o controle e também quando a conexão com o servidor for
perdida. Para isso, será previsto um tempo limite sem comandos válidos.
Esse comportamento deve funcionar independentemente da simulação de ataques.

\section{Metodologia}
O desenvolvimento começará pela organização do protótipo, pelo projeto
do mecanismo de direção acionado por servo motor e pela montagem
do chassi definitivo em impressão 3D. Serão avaliadas a alimentação com
duas baterias 18650, a estabilidade do ESP32 e a resposta da ponte H L298N
aos comandos de sentido de rotação, além do esterçamento produzido pelo servo.

Em seguida, a equipe implementará a comunicação local por Wi-Fi e
WebSocket entre o servidor HIJACK e os carrinhos. O controle terá velocidade
fixa, comandos de frente e ré para a tração e ajuste da direção pelo servo. Será verificado
o desligamento dos motores tanto ao soltar o controle quanto ao perder
a comunicação com o servidor.

Na plataforma web, será implementada a sequência guiada de ataque simulado
com inversões escolhidas individualmente. Depois, serão adicionadas a
troca de senha e as simulações de bloqueio por integridade e de proteção
contra repetição.

Os testes terão uma parte funcional, que verificará o movimento, a
conexão e a resposta às ações do painel, e uma parte quantitativa. Serão
registrados o tempo de resposta de comandos, as tentativas de alteração
bem-sucedidas e bloqueadas, as falhas de conexão e a estabilidade durante
o uso. Como se trata de um pré-projeto, esses resultados ainda serão
obtidos nas etapas de implementação e ensaio.

\section{Resultados Esperados}
Espera-se obter dois carrinhos controlados pela plataforma HIJACK, com
comunicação local estável, tração para frente e para trás e curvas
por esterçamento das rodas.
O sistema deverá permitir observar a inversão simulada de comandos
antes da ativação da defesa e a restauração do comportamento esperado
após o bloqueio da alteração no servidor.

Também se espera que a troca da senha padrão mostre, de forma prática,
a diferença entre uma credencial previsível e uma credencial mais forte.
As medidas de integridade e proteção contra repetição serão identificadas
como demonstrações simuladas, sem atribuir a elas garantias de segurança
que não tenham sido implementadas. Os testes deverão produzir registros
que ajudem a avaliar o funcionamento da comunicação e do controle físico.

\section{Considerações Éticas e de Segurança}
Todos os testes de interferência serão realizados exclusivamente com os
carrinhos e a infraestrutura controlados pelo grupo. A rede local será dedicada à atividade, de modo a evitar interferência
em dispositivos que não fazem parte da demonstração. O propósito é estudar vulnerabilidades em um ambiente
autorizado e mostrar como preveni-las, sem interferir em redes ou
equipamentos de terceiros.

\newpage
\renewcommand{\refname}{Referências}
\begin{thebibliography}{99}
\addcontentsline{toc}{section}{Referências}
\small\sloppy

\bibitem{wired2015jeep}
GREENBERG, A. Hackers remotely kill a Jeep on the highway --- with me in it.
\emph{Wired}, 21 jul.\ 2015. Disponível em:
\url{https://www.wired.com/2015/07/hackers-remotely-kill-jeep-highway/}.

\bibitem{synack2026}
HELP NET SECURITY. AI shrinks vulnerability exploitation window to hours.
18 maio 2026. Disponível em:
\url{https://www.helpnetsecurity.com/2026/05/18/synack-2025-ai-driven-vulnerability-trends-report/}.

\bibitem{thn2026bas}
THE HACKER NEWS. AI broke vulnerability management. That's why CISOs are
moving budget to BAS. jun.\ 2026. Disponível em:
\url{https://thehackernews.com/2026/06/ai-broke-vulnerability-management-thats.html}.

\bibitem{crowdstrike2026}
TECH INSIDER. AI weaponizes bugs in 48 hours, breaches hit \$6.04M. out.\ 2026.
Disponível em:
\url{https://tech-insider.org/ai-cyberattacks-48-hour-exploit-crowdstrike-2026/}.

\bibitem{csa2026praison}
CLOUD SECURITY ALLIANCE. Sub-4-hour weaponization of agentic AI frameworks.
15 maio 2026. Disponível em:
\url{https://labs.cloudsecurityalliance.org/research/csa-research-note-agentic-framework-rapid-exploitation-20260/}.

\bibitem{thn2025gtg1002}
THE HACKER NEWS. Chinese hackers use Anthropic's AI to launch automated cyber
espionage campaign. nov.\ 2025. Disponível em:
\url{https://thehackernews.com/2025/11/chinese-hackers-use-anthropics-ai-to.html}.

\bibitem{bdtech2025}
TECHTALKS. How hackers turned Claude Code into a semi-autonomous cyber-weapon.
nov.\ 2025. Disponível em:
\url{https://bdtechtalks.substack.com/p/how-hackers-turned-claude-code-into}.

\bibitem{aimag2025}
AI MAGAZINE. Anthropic uncovers landmark AI-led cyberattack campaign. nov.\ 2025.
Disponível em:
\url{https://aimagazine.com/news/anthropic-disrupts-first-ai-orchestrated-cyber-espionage}.

\bibitem{checkpoint2026}
HELP NET SECURITY. AI used to help plan the break-in, now it's doing the
break-in. 15 jul.\ 2026. Disponível em:
\url{https://www.helpnetsecurity.com/2026/07/15/check-point-ai-security-report-2026/}.

\bibitem{lmmc2026litellm}
LM MARKET CAP. LiteLLM PyPI supply chain attack: what happened and what to do.
24 mar.\ 2026. Disponível em:
\url{https://lmmarketcap.com/zh/litellm-supply-chain-attack}.

\bibitem{arthur2026litellm}
ARTHUR AI. The LiteLLM supply chain attack: what happened and why it matters.
1 abr.\ 2026. Disponível em:
\url{https://www.arthur.ai/column/litellm-supply-chain-attack-pypi-compromise-2026}.

\bibitem{netenrich2026litellm}
NETENRICH. LiteLLM PyPI supply chain attack: what happened \& how to fix it.
2026. Disponível em:
\url{https://netenrich.com/blog/litellm-pypi-supply-chain-attack}.

\bibitem{hackernoon2026litellm}
HACKERNOON. I spent 48 hours responding to the LiteLLM supply chain attack.
mar.\ 2026. Disponível em:
\url{https://hackernoon.com/i-spent-48-hours-responding-to-the-litellm-supply-chain-attack-here-is-everything-i-know}.

\bibitem{bastion2026litellm}
BASTION. LiteLLM PyPI supply chain attack: what happened and how to protect
your organization. 2026. Disponível em:
\url{https://bastion.tech/blog/litellm-pypi-supply-chain-attack}.

\bibitem{infoq2026litellm}
INFOQ. PyPI supply chain attack compromises LiteLLM, enabling the exfiltration
of sensitive information. mar.\ 2026. Disponível em:
\url{https://infoq.com/news/2026/03/litellm-supply-chain-attack/}.

\bibitem{wiki2026php2shell}
WIKIPEDIA. PHP2Shell. Disponível em:
\url{https://en.wikipedia.org/wiki/PHP2Shell}. Acesso em: 1 nov.\ 2026.

\bibitem{fortune2026php}
FORLINI, E. AI tools discovered the PHP vulnerability that could take down 70\%
of the internet. \emph{Fortune}, 21 out.\ 2026. Disponível em:
\url{https://fortune.com/2026/10/21/ai-tools-discovered-the-php-vulnerability-that-could-take-down-70-of-the-internet-now-the-race-is-on-to-patch-it-before-ai-powered-hackers-attack/}.

\bibitem{wiz2026php2shell}
WIZ RESEARCH. PHP2Shell: AI uncovers critical RCE vulnerabilities affecting 70\%
of the Internet. 21 out.\ 2026. Disponível em:
\url{https://www.wiz.io/blog/php2shell-php-rce}.

\bibitem{techxplore2026php}
TECHXPLORE. AI-assisted discovery exposes decade-old PHP flaw that could put
millions of websites at risk. out.\ 2026. Disponível em:
\url{https://techxplore.com/news/2026-10-ai-discovery-exposes-decade-php.html}.

\bibitem{axios2026php}
AXIOS. The internet's hidden security nightmare --- and how AI exposed it.
22 out.\ 2026. Disponível em:
\url{https://www.axios.com/2026/10/22/php-vulnerability-security-update-websites}.

\bibitem{socprime2026node}
SOC PRIME. CVE-2026-111497: Node.js Brotli decoder flaw enables remote code
execution. 23 out.\ 2026. Disponível em:
\url{https://socprime.com/blog/cve-2026-111497-node-js-brotli-decoder-flaw-enables-remote-code-execution/}.

\bibitem{akamai2026node2shell}
AKAMAI. Inside Node2Shell: a closer look at CVE-2026-111497. out.\ 2026.
Disponível em:
\url{https://www.akamai.com/blog/security-research/inside-node2shell-closer-look-cve-2026-111497}.

\bibitem{explainx2026node}
EXPLAINX. Node.js security updates, October 2026: what was fixed, who should
update, and how. out.\ 2026. Disponível em:
\url{https://explainx.ai/blog/nodejs-security-release-october-2026-brotli-rce-tls-hostname-bypass-emudhra-2026}.

\bibitem{sysdig2026node2shell}
SYSDIG. Node2Shell (CVE-2026-111497): detecting the Node.js Brotli RCE.
28 out.\ 2026. Disponível em:
\url{https://www.sysdig.com/blog/node2shell-cve-2026-111497-detecting-the-node-js-brotli-rce}.

\bibitem{rapid72026node}
RAPID7. CVE-2026-111497: critical out-of-bounds write in the Node.js Brotli
decoder. 26 out.\ 2026. Disponível em:
\url{https://www.rapid7.com/blog/post/etr-cve-2026-111497-critical-out-of-bounds-write-in-the-node-js-brotli-decoder/}.

\bibitem{cloudflare_mirai}
KAN, M. Chinese firm admits its hacked products were behind Friday's massive
DDoS attack. \emph{PCWorld}, out.\ 2016. Disponível em:
\url{https://www.pcworld.com/article/3134039/chinese-firm-admits-its-hacked-products-were-behind-fridays-massive-ddos-attack.html}.

\bibitem{drishti2026}
DRISHTI: AI-led human-directed vulnerability auditing for 5G cores.
\emph{arXiv}, 2026. Disponível em:
\url{https://arxiv.org/pdf/2608.30112}.

\end{thebibliography}

\end{document}